\input{preamble}
\renewcommand{\D}{\mathcal{D}}

\begin{document}

\pagestyle{fancy}
\fancyhead[L]{Seconde}
\fancyhead[C]{\textbf{TD n°7 : Arithmétique}}
\fancyhead[R]{\today}

\exe{}{
Sans utiliser de calculatrice, déterminer la parité des nombres ci-dessous :
\begin{multicols}{2}
\begin{enumerate}[label=(\alph*)]
\item $13~251 + 24~763$
\item $32~280 - 75~813$
\item $128~234 \times 850~437$
\item $(21~756)^3$
\item $(124~843)^2 - 21~033$
\item $4~218 + 5~433 + 7~829 + 2~195$
\end{enumerate}
\end{multicols}
}{exe:1}{
\begin{multicols}{2}
\begin{enumerate}[label=(\alph*)]
\item pair
\item impair
\item impair
\item pair
\item pair
\item impair
\end{enumerate}
\end{multicols}
}

\exe{}{
	On introduit la notation suivante : $\D_n = \set{ k \in \N \tq \text{ k divise n}}$.
	Donner $\D_1, \D_2, \D_3, \D_4, \D_5, \D_6, \D_7, \D_8,$ et $\D_9$.
}{exe:2}{
	\begin{multicols}{2}
	$D_1 = \{ 1 \}$
	
	$D_2 = \{ 1 ; 2\}$
	
	$D_3 = \{ 1 ; 3\}$
	
	$D_4 = \{ 1 ; 4 ; 2 \}$
	
	$D_5 = \{ 1 ; 5\}$
	
	$D_6 = \{ 1 ; 6 ; 2 ; 3\}$
	
	$D_7 = \{ 1 ; 7\}$
	
	$D_8 = \{ 1 ; 8 ; 2 ; 4\}$
	
	$D_9 = \{ 1 ; 3 ; 9\}$
	\end{multicols}
}

\exe{}{
	Décomposer les nombres suivants en produit de facteurs premiers.
	\begin{multicols}{3}
	\begin{enumerate}[label=\roman*)]
		\item  $64 $
		\item $25$
		\item $81$
		\item  $6 \times 121 \times 64$
		\item $45 \times 125 \times 9$
		\item $10^2 \times 15^3$
	\end{enumerate}
	\end{multicols}
}{exe:3}{

	\begin{multicols}{3}
	\begin{enumerate}[label=\roman*)]
		\item  $64 = 2^6 $
		\item $25 = 5^2$
		\item $81 = 9^2$
		\item  $6 \times 121 \times 64 = 2^7 \times 3 \times 11^2$
		\item $45 \times 125 \times 9 = 3^4 \times 5^4$
		\item $10^2 \times 15^3 = 2^2 \times 3^3 \times 5^5$
	\end{enumerate}
	\end{multicols}
}

\exe{}{
On donne :
\[ 3~983 = 43 \times 7 + 14 \times 263 \quad ; \quad 13 \times 257 + 13 \times 48 = 3~978 \quad ; \quad 3~978 \times 3~983 = 15~844~374 \]
Justifier que 3~983 est divisible par 7, que 3~978 est divisible par 13 et que 15~844~374 est divisible par 91.
}{exe:4}{
On fait apparaitre les formes $d \times k$ avec $d$ le diviseur voulu :
\begin{align*} 
3~983 & = 43 \times 7 + 7 \times 2 \times 263 & = 7 \times (43 + 2 \times 263) \\
3~978 & = 13 \times 257 + 13 \times 48 & = 13 \times (257 + 48) \\
3~983 \times 3~978 & = 7 \times (43 + 2 \times 263) \times 13 \times (257 + 48) & = 91 \times (43 + 2 \times 263) \times (257 + 48) 
\end{align*}
car $7 \times 13 = 91$.
}

\exe{}{
	On considère le nombre $n=2^6 \times 3^2 \times 7^7$.
	Quelles sont les affirmations exactes ? Justifier.
	\begin{multicols}{2}
	\begin{enumerate}[label=\alph*)]
		\item $2^3 \times 3^3$ divise $n$
		\item $2^3 \times 7^3$ divise $n$
		\item $2^3 \times 7$ divise $n$
		\item $3^2 \times 7^7$ divise $n$
	\end{enumerate}
	\end{multicols}
}{exe:5}{

	\textbf{Remarque importante} : si $d$ divise $n$, on peut écrire 
		\[ n = d \times k, \]
	pour $k \in \N$ un entier naturel.
	
	En décomposant $d$ et $k$ en produit de nombres premiers, on trouve la décomposition de $n$ (voir exercice \ref{exe:3}).
	Ainsi, si un nombre premier apparaît dans la décomposition de $d$, il doit apparaître au moins autant de fois dans la décomposition de $n$.
	

	\begin{enumerate}[label=\alph*)]
		\item $2^3 \times 3^3$ ne divise pas $n$ car $3$ n'apparaît que $2$ fois dans la décomposition de $n$.
		\item $2^3 \times 7^3$ divise $n$ car $n = (2^3 \times 7^3) \times (2^3 \times 3^2 \times 7^4)$.
		\item $2^3 \times 7$ divise $n$ car $n = (2^3 \times 7) \times ( 2^3 \times 3^2 \times 7^6)$.
		\item $3^2 \times 7^7$ divise $n$ car $n=(3^2 \times 7^7) \times 2^6 $.
	\end{enumerate}
}

\exe{}{
	Exprimer chaque valeur suivante sous forme de fraction irréductible.
	\begin{multicols}{4}
	\begin{enumerate}[label=\alph*), itemsep=10pt]
		\item $1-\dfrac{4}{44}$
		\item $\dfrac{15}{27}$
		\item $\dfrac{8}{3} + 6$
		\item $\dfrac{65}{5^2}$
		\item $\dfrac{6}{32}\times\dfrac{2}{3}$
		\item $3-\dfrac{38}{20}$
		\item $\left(\dfrac{10}{44}\right)^{-1}$
		\item $\dfrac{6}{18}\times2 + 1$
		\item $\left(\dfrac{66}{99}\right)^2$
		\item $\dfrac{23}{18} + \dfrac12$
		\item $\dfrac{34}{20}$
		\item $2 - \dfrac{9}{5}$
	\end{enumerate}
	\end{multicols}
}{exe:6}{
	\begin{multicols}{4}
	\begin{enumerate}[label=\alph*), itemsep=10pt]
		\item $\dfrac{10}{11}$
		\item $\dfrac{5}{9}$
		\item $\dfrac{26}{3}$
		\item $\dfrac{13}{5}$
		\item $\dfrac{1}{18}$
		\item $\dfrac{-11}{10}$
		\item $\dfrac{22}{5}$
		\item $\dfrac{5}{3}$
		\item $\dfrac{4}{9}$
		\item $\dfrac{16}{9}$
		\item $\dfrac{17}{10}$
		\item $\dfrac{1}{5}$
	\end{enumerate}
	\end{multicols}
}

\exe{}{
Soit $p$ un nombre premier. Justifier que s'il existe deux entiers $a$ et $b$ tels que $p=a \cdot b$ alors $a=1$ ou $b=1$.
}{exe:7}{
Par définition d'un nombre premier, si $a$ et $b$ sont différents de $1$, alors $p$ admettrait un autre diviseur que 1 ou lui-même.
}

\exe{, difficulty=1}{
	Soit $n\in\N$ un entier naturel.
	Montrer que $\dfrac{n(n+1)}2$ est toujours un entier.
}{exe:8}{
	Comme $n$ et $n+1$ se suivent, l'un des deux est pair.
	Leur produit est donc pair, et la division par 2 donne un entier.
}


\exe{, difficulty=1}{
Un joueur a besoin d'un dé à 3 faces. Comme il ne possède que des dés à 6 faces, il décide de \og simuler \fg~un dé à 3 faces en suivant la règle suivante :
\begin{itemize}
\item les faces 1 et 2 donnent 1 ;
\item les faces 3 et 4 donnent 2 ;
\item les faces 5 et 6 donnent 3.
\end{itemize}
En suivant la même logique, déterminer les différents dés que vous pouvez simuler avec un dé à 24 faces.
}{exe:9}{
On a $24 = 2^3 times 3$ et $\D_24 = \{1, 2, 3, 4, 6, 8, 12 \}$, on peut donc simuler un dé à : 1 face, 2 faces, 3 faces, 4 faces, 6 faces, 8 faces et 12 faces.
}

\exe{, difficulty=1}{
\begin{enumerate}
\item Donner la décomposition en facteurs premiers de 6 et 15. En déduire l'ensemble des diviseurs \textbf{communs} à 6 et 15.
\item Utiliser cette méthode pour déterminer l'ensemble des diviseurs communs des nombres 420 et 490. Parmi l'ensemble de ces diviseurs, lequel est le plus grand ?
\item Soient $p$ et $q$ deux nombres premiers. Quel est le plus grand diviseur commun de $p$ et $q$ ?
\end{enumerate}
}{exe:10}{
\begin{enumerate}
\item $6 = 2 \times 3$ et $15 = 3 \times 5$, donc l'ensemble des diviseurs communs de 6 et 15 est $\set{1, 3}$.
\item $420 = 2^2 \times 3 \times 5 \times 7$ et $490 = 2 \times 5 \times 7^2$, 
l'ensembles des diviseurs communs est $\set{1, 2, 5, 7, 10, 14, 35, 70}$. Le plus grand est 70.
\item Si $p=q$ alors c'est $p$, sinon c'est 1.
\end{enumerate}
}

\exe{, difficulty=1}{
	On cherche à résoudre l'équation cubique suivante, où $x\in\Z$ est entier :
		\[ x^3 + 3x^2 - x - 108 = 0. \]
	\begin{enumerate}
		\item
		Montrer que l'équation est équivalente à 
			\[  (x-1)(x+1)(x+3) = 105. \]
		\item	
		Donner la décomposition en facteurs premiers de 105.
		\item
		En déduire toutes les valeurs $x\in\Z$ solutions de l'équation originale.
	\end{enumerate}
}{exe:11}{
		\begin{enumerate}
		\item
		On développe le produit $(x-1)(x+1)(x+3)$, on trouve :
		\begin{align*}
			(x-1)(x+1)(x+3) & = (x-1)(x^2+3x+x+3) \\
			& = (x-1)(x^2+4x+3) \\
			& = x^3+4x^2+3x - x^2 - 4x -3 \\
			& = x^3 + 3x^2 - x - 3
		\end{align*}
		Donc $(x-1)(x+1)(x+3) = 105 \iff x^3 + 3x^2 - x - 3 = 105 \iff x^3 + 3x^2 - x - 108 = 0$
		\item	
		$105 = 3 \times 5 \times 7$.
		\item
		On cherche $x \in \Z$ tel que $(x+1)(x-1)(x+3) = 3 \times 5 \times 7$. En tenant compte de la relation d'ordre $x-1 < x+1 < x+3$ et par unicité de la décomposition en facteurs premiers, on en déduit que :
		\begin{align*}
		x-1 & = 3 \\
		x+1 & = 5 \\
		x+3 & = 7 
		\end{align*}
		Soit $x=4$.
	\end{enumerate}
}

\exe{, difficulty=2}{
Soit $k \in \N$. Déterminer une règle générale pour connaître la parité de la somme de $k$ nombres impairs quelconques.
}{exe:12}{
Si $k$ est pair, on peut regrouper les impairs par couples, la somme des deux impairs du coupe est un nombre pair. On somme ensuite les $\frac{k}{2}$ nombres pairs obtenus et on trouve un nombre pair.
Si $k$ est impair, on aura $\frac{k-1}{2}$ couples dont la somme est paire et un impair isolé, donc la somme totale sera impaire.
}

\exe{,difficulty=2}{
On donne $1820 = 2^2 \times 5 \times 7 \times 13$.
\begin{enumerate}
\item Donner l'ensemble des nombres premiers divisant 1820. Expliquer.
\item On considère un entier $n$ dont la décomposition en facteurs premiers comporte les nombres premiers $p_1, p_2, \dots, p_d$.
Est-il possible qu'un nombre premier $p$ n'appartenant pas à la liste $(p_1, p_2, \dots, p_d)$ divise $n$ ? Justifier.
\item En déduire le lemme d'Euclide : \textit{Soient $a$ et $b$ deux entiers. Si un nombre premier $p$ divise le produit $ab$, alors $p$ divise $a$ ou $b$}.
\end{enumerate}
}{exe:13}{
\begin{enumerate}
\item Il suffit d'isoler les nombres premiers dans la décomposition en facteurs premiers de 1820. Ainsi 2, 5, 7 et 13 divisent 1820.
\item On peut écrire $n = p_1^{\alpha_1} \cdot p_2^{\alpha_2} \cdot \dots \cdot p_d^{\alpha_d}$ avec les nombres $\alpha_1, \dots, \alpha_d$ des entiers supérieurs ou égaux à 1.
Si $p$ divise $n$, on peut écrire $n=p \cdot k$, avec $k \in \Z$. En d'autres termes :
\[ p \cdot k = p_1^{\alpha_1} \cdot p_2^{\alpha_2} \cdot \dots \cdot p_d^{\alpha_d} \]
et en divisant, par exemple, à gauche et à droite par $p_1$ on conclut que $p$ est divisible par $p_1$ ce qui est absurde. On aboutit donc à une contradiction.
\item Si $p$ divise le produit $ab$, selon la question précédente, il divise un des nombres premiers apparaissant dans l'une ou l'autre des décomposition en facteurs premiers de $a$ et $b$.
Or, le décomposition en facteurs premiers de $ab$ ne peut pas comporter d'autres nombres premiers que ceux apparaissant dans celles de $a$ et $b$ (par définition d'un nombre premier et par unicité de cette décomposition).
\end{enumerate}
}


%%%%%%%%%%%%

\newpage
\fancyhead[C]{\textbf{Solutions}}
\shipoutAnswer

\end{document}
